\documentclass[12pt,letterpaper]{article}
\usepackage{graphicx,textcomp}
\usepackage{natbib}
\usepackage{setspace}
\usepackage{fullpage}
\usepackage{color}
\usepackage[reqno]{amsmath}
\usepackage{amsthm}
\usepackage{fancyvrb}
\usepackage{amssymb,enumerate}
\usepackage[all]{xy}
\usepackage{endnotes}
\usepackage{lscape}
\newtheorem{com}{Comment}
\usepackage{float}
\usepackage{hyperref}
\newtheorem{lem} {Lemma}
\newtheorem{prop}{Proposition}
\newtheorem{thm}{Theorem}
\newtheorem{defn}{Definition}
\newtheorem{cor}{Corollary}
\newtheorem{obs}{Observation}
\usepackage[compact]{titlesec}
\usepackage{dcolumn}
\usepackage{tikz}
\usetikzlibrary{arrows}
\usepackage{multirow}
\usepackage{xcolor}
\newcolumntype{.}{D{.}{.}{-1}}
\newcolumntype{d}[1]{D{.}{.}{#1}}
\definecolor{light-gray}{gray}{0.65}
\usepackage{url}
\usepackage{listings}
\usepackage{color}

\definecolor{codegreen}{rgb}{0,0.6,0}
\definecolor{codegray}{rgb}{0.5,0.5,0.5}
\definecolor{codepurple}{rgb}{0.58,0,0.82}
\definecolor{backcolour}{rgb}{0.95,0.95,0.92}

\lstdefinestyle{mystyle}{
	backgroundcolor=\color{backcolour},   
	commentstyle=\color{codegreen},
	keywordstyle=\color{magenta},
	numberstyle=\tiny\color{codegray},
	stringstyle=\color{codepurple},
	basicstyle=\footnotesize,
	breakatwhitespace=false,         
	breaklines=true,                 
	captionpos=b,                    
	keepspaces=true,                 
	numbers=left,                    
	numbersep=5pt,                  
	showspaces=false,                
	showstringspaces=false,
	showtabs=false,                  
	tabsize=2
}
\lstset{style=mystyle}
\newcommand{\Sref}[1]{Section~\ref{#1}}
\newtheorem{hyp}{Hypothesis}


\title{Problem Set 4}
\date{Due: November 26, 2026}
\author{Applied Stats/Quant Methods 1}


\begin{document}
	\maketitle
	\section*{Instructions}
	\begin{itemize}
		\item Please show your work! You may lose points by simply writing in the answer. If the problem requires you to execute commands in \texttt{R}, please include the code you used to get your answers. Please also include the \texttt{.R} file that contains your code. If you are not sure if work needs to be shown for a particular problem, please ask.
		\item Your homework should be submitted electronically on GitHub.
		\item This problem set is due before 23:59 on Thursday November 26, 2026. No late assignments will be accepted.
	\end{itemize}



	\vspace{.5cm}
\section*{Question 1: Economics}
\vspace{.25cm}
\noindent 	
In this question, use the \texttt{prestige} dataset in the \texttt{car} library. First, run the following commands:

\begin{verbatim}
install.packages(car)
library(car)
data(Prestige)
help(Prestige)
\end{verbatim} 


\noindent We would like to study whether individuals with higher levels of income have more prestigious jobs. Moreover, we would like to study whether professionals have more prestigious jobs than blue and white collar workers.

\newpage
\begin{enumerate}
	
	\lstinputlisting[language=R, firstline=46, lastline=53]{PS04_PT.R}	
	
	\item [(a)]
	Create a new variable \texttt{professional} by recoding the variable \texttt{type} so that professionals are coded as $1$, and blue and white collar workers are coded as $0$ (Hint: \texttt{ifelse}).
	
	\lstinputlisting[language=R, firstline=63, lastline=63]{PS04_PT.R}
	\begin{verbatim}
	bc prof   wc  NAs 
	44   31   23    4 
	\end{verbatim}

	\lstinputlisting[language=R, firstline=58, lastline=61]{PS04_PT.R}
	
	\begin{verbatim}
		      0  1
		bc   44  0
		prof  0 31
		wc   23  0
	\end{verbatim}
	
	\textbf{We can see from our data that we have 67 blue collar and white collar workers, who are all coded as 0, and 31 professional workers, who are all coded as 1. We also know from our summary of the data that there are 4 NA's in the data, which are omitted from this table.}
	
	\item [(b)]
	Run a linear model with \texttt{prestige} as an outcome and \texttt{income}, \texttt{professional}, and the interaction of the two as predictors (Note: this is a continuous $\times$ dummy interaction.)
	
	
	\lstinputlisting[language=R, firstline=68, lastline=69]{PS04_PT.R}
	\begin{verbatim}
		Call:
		lm(formula = prestige ~ income + professional + income * professional, 
		data = Prestige)
		
		Residuals:
		Min      1Q  Median      3Q     Max 
		-14.852  -5.332  -1.272   4.658  29.932 
		
		Coefficients:
		Estimate Std. Error t value           Pr(>|t|)    
		(Intercept)         21.1422589  2.8044261   7.539 0.0000000000292686 ***
		income               0.0031709  0.0004993   6.351 0.0000000075482422 ***
		professional        37.7812800  4.2482744   8.893 0.0000000000000414 ***
		income:professional -0.0023257  0.0005675  -4.098 0.0000882872162594 ***
		---
		Signif. codes:  0 ‘***’ 0.001 ‘**’ 0.01 ‘*’ 0.05 ‘.’ 0.1 ‘ ’ 1
		
		Residual standard error: 8.012 on 94 degrees of freedom
		(4 observations deleted due to missingness)
		Multiple R-squared:  0.7872,	Adjusted R-squared:  0.7804 
		F-statistic: 115.9 on 3 and 94 DF,  p-value: < 0.00000000000000022
	\end{verbatim}	
\textbf{The R squared value tells us that the model explains 78.72 percent of the variation in prestige.}

	\item [(c)]
	Write the prediction equation based on the result.
	
	\textbf{Predicted prestige = Intercept + (Income coefficient x Income) + (Professional coefficient x professional) + (Interaction coefficient of professional and income x income x professional)}
	
	\textbf{With our coefficient numbers:}
	
	\textbf{predicted prestige = 21.142 + (0.00317 x income) + (37.781 x professional) + (-0.00233 x income x professional)}
	
	\lstinputlisting[language=R, firstline=75, lastline=81]{PS04_PT.R}
		
	
\newpage
	\item [(d)]
	Interpret the coefficient for \texttt{income}.
	
	\textbf{For blue collar and white collar occupations each additional dollar of income is associated with about a 0.00317 unit increase in prestige on average. Because the p-value is far below 0.05 we reject the null hypothesis that income has no relationship with prestige for blue and white collar occupations.}

	\item [(e)]
	Interpret the coefficient for \texttt{professional}.
	
	\textbf{When income is 0, professional occupations are associated with 37.781 more units of prestige than blue collar and white collar workers on average. Because the p-value is far below 0.05 we reject the null hypothesis that there is no difference in prestige between blue collar/white collar workers and professionals when income is 0.}
	
	\item [(f)]
	What is the effect of a \$1,000 increase in income on prestige score for professional occupations? In other words, we are interested in the marginal effect of income when the variable \texttt{professional} takes the value of $1$. Calculate the change in $\hat{y}$ associated with a \$1,000 increase in income based on your answer for (c).
	
	\textbf{we take our prediction equation:}
	\textbf{predicted prestige = 21.142 + (0.00317 x income) + (37.781 x professional) + (-0.00233 x income x professional)}
	\textbf{and by hand we set income to 1,000, set professional to 1, check difference when income is 0}
	\lstinputlisting[language=R, firstline=104, lastline=107]{PS04_PT.R}
	\begin{verbatim}
		[1] 0.84
	\end{verbatim}
	\textbf{we can also take our function from 1c and check differences in income of 1000}
	\lstinputlisting[language=R, firstline=100, lastline=100]{PS04_PT.R}
	\begin{verbatim}
		[1] 0.8452
	\end{verbatim}
	\textbf{we can do it again with different numbers just to make sure that it works:}
	\lstinputlisting[language=R, firstline=101, lastline=101]{PS04_PT.R}
	\begin{verbatim}
		[1] 0.8452
	\end{verbatim}
	
	\textbf{For professional occupations we can see that a 1000 dollar increase in income is associated with an increase in 0.8452 units of predicted prestige}
	
	\item [(g)]
	What is the effect of changing one's occupations from non-professional to professional when her income is \$6,000? We are interested in the marginal effect of professional jobs when the variable \texttt{income} takes the value of $6,000$. Calculate the change in $\hat{y}$ based on your answer for (c).
	
		\textbf{we take our prediction equation again to do this by hand,  we set income to 6,000 in both equations. We set professional to 1 on one side, and subtract that when professional is set to 0 on the other side.}
	\lstinputlisting[language=R, firstline=117, lastline=120]{PS04_PT.R}
	\begin{verbatim}
[1] 23.801
	\end{verbatim}
	\textbf{We can check this by using our function from earlier and doing the same thing, and we will get an even more precise result:}
		\lstinputlisting[language=R, firstline=116, lastline=116]{PS04_PT.R}
	\begin{verbatim}
[1] 23.82703
	\end{verbatim}
	\textbf{When both professional and non-professional workers have an income of 6000 dollars, professional workers are expected to have 23.82703 higher predicted points of prestige.}
\end{enumerate}

\newpage

\section*{Question 2: Political Science}
\vspace{.25cm}
\noindent 	Researchers are interested in learning the effect of all of those yard signs on voting preferences.\footnote{Donald P. Green, Jonathan	S. Krasno, Alexander Coppock, Benjamin D. Farrer,	Brandon Lenoir, Joshua N. Zingher. 2016. ``The effects of lawn signs on vote outcomes: Results from four randomized field experiments.'' Electoral Studies 41: 143-150. } Working with a campaign in Fairfax County, Virginia, 131 precincts were randomly divided into a treatment and control group. In 30 precincts, signs were posted around the precinct that read, ``For Sale: Terry McAuliffe. Don't Sellout Virgina on November 5.'' \\

Below is the result of a regression with two variables and a constant.  The dependent variable is the proportion of the vote that went to McAuliff's opponent Ken Cuccinelli. The first variable indicates whether a precinct was randomly assigned to have the sign against McAuliffe posted. The second variable indicates
a precinct that was adjacent to a precinct in the treatment group (since people in those precincts might be exposed to the signs).  \\

\vspace{.5cm}
\begin{table}[!htbp]
	\centering 
	\textbf{Impact of lawn signs on vote share}\\
	\begin{tabular}{@{\extracolsep{5pt}}lccc} 
		\\[-1.8ex] 
		\hline \\[-1.8ex]
		Precinct assigned lawn signs  (n=30)  & 0.042\\
		& (0.016) \\
		Precinct adjacent to lawn signs (n=76) & 0.042 \\
		&  (0.013) \\
		Constant  & 0.302\\
		& (0.011)
		\\
		\hline \\
	\end{tabular}\\
	\footnotesize{\textit{Notes:} $R^2$=0.094, N=131}
\end{table}

\vspace{.5cm}
\begin{enumerate}
	\item [(a)] Use the results from a linear regression to determine whether having these yard signs in a precinct affects vote share (e.g., conduct a hypothesis test with $\alpha = .05$).
	
	\textbf{We need to get the t-stats first to then figure out the p-values and see if our results are statistically significant.}
	\textbf{Our tstat equation is: TS = (beta of yard signs in precincts - 0)/SE of yard signs in precincts}
		\lstinputlisting[language=R, firstline=127, lastline=128]{PS04_PT.R}
	\begin{verbatim}
	[1] 2.625
	\end{verbatim}
	\textbf{then to find the 2-sided p-values we do pvalues = 2*pt(abs(TS), n-3, lower.tail = F)}
	\lstinputlisting[language=R, firstline=130, lastline=131]{PS04_PT.R}
		\begin{verbatim}
		[1] 0.00972002
		\end{verbatim}
	\textbf{because our pvalue of 0.0097 is less than 0.05 we can reject the NULL hypothesis that yard signs do not affect vote share.}	
	\item [(b)]  Use the results to determine whether being
	next to precincts with these yard signs affects vote
	share (e.g., conduct a hypothesis test with $\alpha = .05$).
	
	\textbf{first for the t-stat do TS = (beta of signs in adjacent yards - 0)/SE signs in adjacent yards}
	\lstinputlisting[language=R, firstline=136, lastline=137]{PS04_PT.R}
			\begin{verbatim}
	[1] 3.230769
	\end{verbatim}
	
	\textbf{then, again, for the 2-sided p-values do pvalues = 2*pt(abs(TS), n-3, lower.tail = F)}
\lstinputlisting[language=R, firstline=139, lastline=140]{PS04_PT.R}
		\begin{verbatim}
[1] 0.00156946
\end{verbatim}
	\textbf{because our pvalue of 0.0016 is less than 0.05 we can reject the NULL hypothesis that signs in adjacent yards do not affect vote share.}
	\item [(c)] Interpret the coefficient for the constant term substantively.

		\textbf{In precincts that were neither assigned lawn signs nor were they adjacent to lawn signs: Cuccinelli was predicted to receive 0.302 of the vote share. This is the group that the sign effects were compared against.}
	
	
	\item [(d)] Evaluate the model fit for this regression.  What does this	tell us about the importance of yard signs versus other factors that are not modeled?
	
	\textbf{The R squared value is 0.094 which tells us that the lawn sign vairables explain about 9.4 percent variation in the vote share measured across 131 precincts. The rest of the variation in vote share is due to other variables that are not included in the model. While the signs have statistically significant effects they do not explain the entire difference in vote share. We know that the effect of the signs is unlikely to be zero, and the R squared value tells us how much of the variation is explained by our variable. While there are likely other factors such as political partisanship which have a far higher correlation to vote share, this does not mean that the signs are an unimportant variable as 4 percentage points can be very significant in a close election.}
	
\end{enumerate}  


\end{document}
